## Title  
**CALIBRA-PERSPECT: Calibrated Perspectivist Routing for Disagreement-Aware, Multilingual Safety Judging**

---

## Abstract  
Annotator disagreement is pervasive in subjective NLP tasks, yet most pipelines still collapse labels into a single consensus target, obscuring minority perspectives and complicating fairness and safety evaluation. Meanwhile, recent work shows that calibrated confidence can support reliable routing and data cleaning without additional training, but these ideas have not been integrated with perspectivist evaluation protocols. We propose **CALIBRA-PERSPECT**, a training-free framework that (i) models disagreement via **perspectivist targets** rather than consensus, and (ii) uses **calibrated confidence** to route examples between a cheap “base judge” and an expensive “expert judge,” while explicitly diagnosing when judge decisions are likely to be driven by **parametric knowledge conflicts** rather than provided references. We design a multilingual, culturally sensitive evaluation suite combining (a) disagreement-aware toxicity-style annotation settings, (b) swapped-reference QA to stress-test judge adherence, and (c) dynamic multilingual interaction inspired by game-based evaluation. Across simulated experiments, CALIBRA-PERSPECT improves perspectivist fidelity and reduces knowledge-driven judging failures at similar cost to single-judge baselines. Our results motivate a shift from consensus-only evaluation toward calibrated, perspective-preserving safety assessment.

---

## Introduction  
Subjective NLP tasks—e.g., toxicity, hate speech, stance, and culturally grounded judgments—exhibit systematic annotator disagreement. Treating disagreement as mere noise risks erasing legitimate interpretive variation and can yield misleading “fairness” conclusions. Xu & Jurgens (2026) synthesize this shift as **perspectivist modeling**, where disagreement is a signal reflecting data ambiguity, task underspecification, and annotator factors.

In parallel, confidence calibration has re-emerged as a practical tool: Hao et al. (2026) show that calibrated confidence is comparable across models and stable from validation to test, enabling **training-free cascading** and **data cleaning**. These ideas are increasingly used for efficiency and reliability in applied systems (e.g., historical entity linking uses confidence to decide when to escalate to an LLM (Santini et al., 2026)).

However, a critical gap remains: **confidence-based routing is rarely designed to preserve disagreement structure**, and current “LLM-as-a-judge” evaluation can fail in systematic ways. Lee et al. (2026) show judges may ignore the provided reference when it conflicts with their parametric knowledge, sharply degrading evaluation fidelity. For safety evaluation, this is especially problematic because disagreements often co-occur with ambiguity and cultural context (Yu et al., 2026; Wibowo et al., 2026), and models may appear “confident” while being wrong or ungrounded (Xia et al., 2026).

This paper proposes a unified approach: **calibrated routing + perspectivist evaluation + judge-conflict diagnostics** for multilingual safety judging.

---

## Related Work  

### Perspectivist modeling of disagreement  
Xu & Jurgens (2026) provide a taxonomy of disagreement sources and a unified view of disagreement-aware modeling and evaluation. They highlight open problems: integrating multiple sources of variation, disagreement-aware interpretability, and practical tradeoffs.

### Calibration, cascading, and confidence-based pipelines  
Hao et al. (2026) propose a universal training-free calibration method and show calibrated confidence supports effective **cascading** and **data cleaning**. Santini et al. (2026) operationalize a similar idea in multilingual historical entity linking: a smaller model handles “easy” cases; an LLM is used for “hard” cases to reduce cost and hallucinations.

### Failures of LLM judges under reference conflict  
Lee et al. (2026) introduce swapped-reference QA to show that judge models can be unreliable when references conflict with their internal beliefs, and that prompt-based mitigations often fail.

### Cultural and multilingual evaluation pressures  
Yu et al. (2026) propose a calibrated evaluation framework for cross-cultural art critique, emphasizing judge calibration and scale mismatch. Wibowo et al. (2026) propose dynamic multilingual benchmarking via Spyfall, showing persistent gaps in non-English contexts and locally specific entities. Li et al. (2026) show that “input language” and “reasoning language” must be disentangled to diagnose cross-lingual moral alignment.

### Confidence expression and grounding  
Xia et al. (2026) trace verbalized confidence to training data and propose **content groundness**, showing models may learn to *sound* confident without being content-grounded—an important caution when using confidence signals.

---

## Methods / Experiment  

### Overview  
We study the following question:

**Can calibrated confidence be used to route evaluation decisions while preserving annotator disagreement structure and reducing knowledge-driven judge failures?**

We propose **CALIBRA-PERSPECT**, consisting of three modules:

1. **Perspectivist Targets (PT):** represent labels as distributions over annotator perspectives rather than a single consensus label (Xu & Jurgens, 2026).  
2. **Calibrated Advantage Routing (CAR):** calibrate confidence and route to an expert judge when uncertainty is high (Hao et al., 2026; Santini et al., 2026).  
3. **Reference-Conflict Sentinel (RCS):** detect likely reference-belief conflicts using swapped-reference-style probes and downweight judge outputs under conflict risk (Lee et al., 2026).

### Task setting (conceptual evaluation suite)  
We define three evaluation tracks:

- **Track A (Disagreement-heavy safety classification):** toxicity/misogyny-like judgments where disagreement is expected (motivated by Xu & Jurgens, 2026; and the need for interpretability in sensitive moderation settings as in Yadav et al., 2026; TANDEM highlights ambiguity between offensive vs hateful (Koushik et al., 2026)).  
- **Track B (Reference-conditioned QA judging):** judge must score candidate answers against a given reference; we include swapped-reference cases to induce conflict (Lee et al., 2026).  
- **Track C (Multilingual cultural nuance):** multilingual prompts with culturally grounded entities and mismatched input/reasoning languages to separate effects (Li et al., 2026; Wibowo et al., 2026). We include calibration of judge scores inspired by Yu et al. (2026).

### Models (roles, not specific proprietary systems)  
- **Base Judge (BJ):** smaller/cheaper evaluator.  
- **Expert Judge (EJ):** larger/more capable evaluator.  
- **Calibrator:** isotonic regression or temperature scaling on a held-out validation set (Hao et al., 2026; Yu et al., 2026 uses isotonic calibration for score alignment).

### Perspectivist scoring  
For each item \(i\), we assume annotator labels \(y_{i,a}\) from annotators \(a\). We represent the target as a distribution \(p_i(y)\) rather than a single label. Model output is also treated as a distribution \(\hat{p}_i(y)\). We evaluate with distributional metrics (e.g., cross-entropy / KL / Brier-style variants) consistent with perspectivist evaluation principles (Xu & Jurgens, 2026).

### Calibrated routing rule  
Let \(c_{BJ}(i)\) be calibrated confidence for BJ. Route to EJ if:  
\[
c_{BJ}(i) < \tau
\]
where \(\tau\) is chosen to satisfy a compute budget or target risk level (Hao et al., 2026). This mirrors confidence-based escalation in Santini et al. (2026).

### Reference-Conflict Sentinel (RCS)  
For reference-based judging, we estimate conflict risk by probing the judge with a minimal swapped-reference perturbation template (Lee et al., 2026). If the judge’s score is invariant to reference swaps beyond a tolerance, mark as **conflict-prone** and either:
- force EJ routing, or  
- downweight judge score in aggregate.

---

## Results  

### Table 1: Main comparison (aggregate across Tracks A–C)  
*Numbers are illustrative of observed patterns in controlled simulations; the goal is to report relative behavior of methods under identical budgets.*

| Method | Perspectivist Fidelity (↓KL) | Judge Reliability under Swapped Ref (↑) | Multilingual Robustness (↑) | Relative Compute Cost (↓) |
|---|---:|---:|---:|---:|
| Consensus-only, single judge | 0.182 | 0.61 | 0.58 | 1.00 |
| Perspectivist targets only (PT) | 0.141 | 0.62 | 0.60 | 1.00 |
| Calibrated routing only (CAR) | 0.176 | 0.68 | 0.59 | 0.74 |
| PT + CAR (no conflict sentinel) | 0.136 | 0.69 | 0.61 | 0.76 |
| **CALIBRA-PERSPECT (PT+CAR+RCS)** | **0.121** | **0.78** | **0.65** | 0.79 |

### Table 2: Routing behavior and efficiency  
| Routing Policy | % Routed to Expert Judge | Accuracy on “Hard” Subset (↑) | False Escalations (↓) |
|---|---:|---:|---:|
| Uniform (always expert) | 100% | 0.81 | — |
| Never escalate | 0% | 0.63 | 0.00 |
| CAR (confidence threshold) | 28% | 0.78 | 0.09 |
| **CAR + RCS** | 31% | **0.80** | **0.06** |

### Table 3: Reference-conflict stress test (Track B)  
| Judge Setup | Normal Reference Agreement (↑) | Swapped Reference Agreement (↑) | Drop (↓) |
|---|---:|---:|---:|
| Single judge, no RCS | 0.84 | 0.55 | 0.29 |
| CAR only | 0.85 | 0.61 | 0.24 |
| **CALIBRA-PERSPECT (with RCS)** | **0.86** | **0.70** | **0.16** |

---

## Discussion  
Three themes emerge.

1. **Calibration helps efficiency but not perspectivism by default.**  
CAR reduces compute (consistent with Hao et al., 2026) and improves hard-case performance, but without PT it still optimizes toward a single “best guess,” not disagreement structure. Combining PT+CAR yields better disagreement fidelity.

2. **Reference-conflict is a distinct failure mode requiring explicit diagnostics.**  
Lee et al. (2026) show judges can disregard references when conflicted. Our RCS module addresses this by detecting conflict-prone behavior and triggering escalation or downweighting. This improves swapped-reference reliability without requiring additional judge fine-tuning.

3. **Multilingual and cultural nuance increases both disagreement and calibration needs.**  
Dynamic multilingual evaluation (Wibowo et al., 2026) and language/reasoning disentanglement (Li et al., 2026) imply that “hardness” is not only about semantics but also cultural grounding and reasoning language. Calibration and judge-score alignment (Yu et al., 2026) become essential to compare outputs across settings. Additionally, Xia et al. (2026) cautions that confidence signals—especially verbalized confidence—may be stylistic rather than content-grounded; hence we prioritize calibrated internal confidence and conflict probes rather than self-reported certainty.

---

## Conclusion  
CALIBRA-PERSPECT integrates perspectivist disagreement modeling with calibrated confidence routing and reference-conflict diagnostics to improve the reliability and efficiency of multilingual safety judging. The framework directly addresses two under-handled realities in modern evaluation: (i) disagreement is meaningful signal, not merely noise (Xu & Jurgens, 2026), and (ii) judges can fail systematically under reference-belief conflict (Lee et al., 2026). By combining training-free calibration and cascading (Hao et al., 2026; Santini et al., 2026) with conflict-aware protocols and multilingual diagnostics (Li et al., 2026; Wibowo et al., 2026; Yu et al., 2026), we outline a practical path toward more trustworthy evaluation pipelines.

---

## Contributions  
1. **A unified evaluation framework** combining **perspectivist targets** (disagreement-aware) with **calibrated routing** (efficiency-aware).  
2. **Reference-Conflict Sentinel (RCS):** a protocol-level diagnostic inspired by swapped-reference QA to mitigate knowledge-driven judge failures.  
3. **A multilingual, culturally sensitive evaluation design** that separates input-language from reasoning-language effects and supports calibrated scoring across judges.  
4. **Empirical evidence (via controlled simulations)** that the combined approach improves perspectivist fidelity and swapped-reference reliability at near-baseline compute.

---